\section{Reproducing the numbers}\label{app:reproducibility}

Every table and figure is generated by a script in \code{paper/scripts/}, which writes a CSV file
to \code{paper/data/}. The figures and some tables read those files directly. The scripts need
Python~3 with NumPy and SciPy (\code{scripts/requirements.txt}). \code{levelB\_detector.py} also
needs a Rust toolchain: it builds \code{scripts/dif-harness}, a small program that runs the
production detector \code{scoring::dif::mixture\_dif} and the protocol's verdict rule on
generated batches. All randomness is seeded.

\begin{center}\small
\begin{tabular}{@{}lll@{}}
\toprule
Script & Produces & Runtime\\
\midrule
\code{levelA\_relativity.py} & identities of \cref{lem:gauge}; \cref{tab:relativity} & $<1$ min\\
\code{levelA\_leak.py} & \cref{tab:leak}, \cref{fig:leak} & $\approx8$ min\\
\code{levelB\_information.py} & \cref{tab:information} & seconds\\
\code{levelB\_proxy.py} & \cref{tab:proxy} & $\approx2$ min\\
\code{levelB\_detector.py} & \cref{tab:anchors}; production-detector counts in the text & $\approx8$ min\\
\code{levelB\_differential.py} & differential gap (\cref{sec:proxy}) & $<1$ min\\
\code{levelC\_bss.py} & \cref{fig:bss}, \cref{tab:bss} & $\approx1$ min\\
\code{assignment.py} & \cref{tab:assignment} & seconds\\
\bottomrule
\end{tabular}
\end{center}

\code{common.py} re-implements the bridging objective and the simulation's two-class DIF model in
NumPy with analytic gradients and tight tolerances, so that stationarity conditions can be checked
to optimizer precision. Its function \code{sim\_levelA\_dataset()} rebuilds, from the same
random stream, the dataset of the repository simulation \code{sim/bridging\_irt\_dif.py}. The paper is built with
\code{latexmk -pdf main.tex} in \code{paper/}.
